\begin{frame}[fragile]{Comments (Waarvoor?)}
    Comments zijn nuttig voor twee belangrijke redenen:
    \begin{itemize}[label={\textbullet}]
        \pause 
        \item Tekst \textbf{opslaan};
        \pause 
        \begin{itemize}[label=$\blacktriangleright$]
            \item  Als je een stuk van een opdracht hebt uitgetypt, en je komt erachter dat het niet helemaal klopt, 
                dan is het vaak handig om dit niet allemaal te verwijderen, maar het gewoon in een comment te 
                zetten. Zo blijft het toch opgeslagen in je bestand, zonder dat het weergegeven wordt.
        \end{itemize}
        \pause 
        \item Om \textbf{foutmeldingen} te verhelpen;
        \pause
        \begin{itemize}[label=$\blacktriangleright$]
            \item Bij een foutmelding die je niet snapt \textrightarrow comment je code \emph{helemaal} 
                en compile je document. Als je nu geen errors krijgt, dan heb je het stuk met de fout te pakken. 
                Je kunt nu langzaam regels gaan uncommenten en je bestand weer compilen om zo te vinden 
                waar de error zit. 
        \end{itemize}
    \end{itemize}
\end{frame}
